\begin{center}
\LARGE
Symbolic Computation of Linear and Nonlinear Modified
                (Partial Differential) Equations

\end{center}

\begin{center}
\Large
Jean-Antoine D\'esid\'eri, Margarita Spiridonova, 

\end{center}

\noindent
Date:  July 18th (Thursday)\\
Time:  11:35-12:00\\

\begin{center}
\large
Abstract
\end{center}

The so-called ``equivalent'' (Lerat, Peyret, 1974) or ``modified'' (Warming, Hyett, 1974) partial
differential equation is a very useful, classical and basic tool of analysis of finite-difference
schemes for time-dependent problems, mostly but not only, of hyperbolic type. Although limited
in application to simple (usually scalar) model equations, it allows in particular a rigorous analysis
of the truncation error, in which dispersive and diffusive errors are identified. Thus new schemes
are systematically compared with well-known ones by this analysis applied to a test equation
(often the wave or heat equation). Furthermore, the approach can be used constructively to build
new schemes in which particular types of error are absent (to a certain order of accuracy). 

Originally, Warming and Hyett developed codes of symbolic computation of modified equations
using, to our knowledge, the language MACSYMA. We have developed a collection of programs
of the same nature for symbolic derivation of the modified equation associated with a given
finite-difference equation defined in a rather general format using the computer algebra system
MAPLE. 

This contribution intends to report on the implementation of these programs and their
experimentation related to various finite-difference schemes. For the case of linear (hyperbolic or
parabolic) partial differential equations, we have examined classical schemes and less classical
examples in which equations in one or two space dimensions including a source term are
approximated by upwind schemes. For the case of nonlinear equations, the MacCormack scheme
applied to Burgers equation was studied as well as Runge-Kutta-type methods applied to general
hyperbolic equations. 

